\chapter{Programs — Index}
%%% generated by the blueprint pipeline from TCSlib.Programs.Util.Index
\section{Overview}

Declarations of the Lean module \texttt{TCSlib.Programs.Util.Index}.

\section{Declarations}

\begin{theorem}[A product and sum bound]
\lean{Nat.mul_add_lt_mul}
\label{Nat.mul_add_lt_mul}
\leanok
\difficulty{1}
For natural numbers $a,b,m,n$, if $a<m$ and $b<n$, then $an+b<mn$.
\end{theorem}

\begin{theorem}[Bound on a product plus a summand]
\lean{Nat.mul_add_le_mul}
\label{Nat.mul_add_le_mul}
\leanok
\difficulty{4}
For natural numbers $a,b,m,n$, if $a<m$ and $b\le n$, then $an+b\le mn$.
\end{theorem}

\begin{theorem}[Quotient of a multiple plus a small remainder]
\lean{Nat.mul_add_div_of_lt}
\label{Nat.mul_add_div_of_lt}
\leanok
\difficulty{3}
For natural numbers $a,b,n$ with $b<n$, the quotient $(an+b)/n$ equals $a$.
\end{theorem}

\begin{theorem}[Uniqueness of quotient and remainder]
\lean{Nat.mul_add_inj_of_lt}
\label{Nat.mul_add_inj_of_lt}
\leanok
\uses{Nat.mul_add_div_of_lt}
\difficulty{3}
For natural numbers $a,b,a',b',n$ with $b<n$ and $b'<n$, if $an+b=a'n+b'$, then $a=a'$
and $b=b'$.
\end{theorem}

\begin{theorem}[Quotient bound from a product bound]
\lean{Nat.div_lt_of_lt_mul'}
\label{Nat.div_lt_of_lt_mul'}
\leanok
\difficulty{1}
For natural numbers $t,m,n$, if $t<mn$, then the quotient $t/n$ is less than $m$.
\end{theorem}

\begin{theorem}[Remainder bound from a product bound]
\lean{Nat.mod_lt_of_lt_mul}
\label{Nat.mod_lt_of_lt_mul}
\leanok
\difficulty{1}
For natural numbers $t,m,n$, if $t<mn$, then the remainder $t\bmod n$ is less than $n$.
\end{theorem}

\begin{theorem}[Bounded quotient and remainder representation]
\lean{Nat.exists_eq_mul_add_of_lt_mul}
\label{Nat.exists_eq_mul_add_of_lt_mul}
\leanok
\uses{Nat.div_lt_of_lt_mul', Nat.mod_lt_of_lt_mul}
\difficulty{1}
For natural numbers $t,m,n$, if $t<mn$, then there exist natural numbers $a<m$ and $b<n$
such that $t=an+b$.
\end{theorem}

\begin{theorem}[Characterization of a quotient and remainder]
\lean{Nat.eq_mul_succ_iff}
\label{Nat.eq_mul_succ_iff}
\leanok
\uses{Nat.mul_add_div_of_lt}
\difficulty{3}
For natural numbers $n,i,t$ with $i<n$, the equality $t=i(n+1)$ holds if and only if
both $t/n=i$ and $t\bmod n=i$.
\end{theorem}

\begin{theorem}[Quotient and remainder below a boundary]
\lean{Nat.succ_div_mod_of_lt}
\label{Nat.succ_div_mod_of_lt}
\leanok
\uses{Nat.mul_add_div_of_lt}
\difficulty{3}
For natural numbers $n$ and $i$, if $i\bmod n+1<n$, then $(i+1)/n=i/n$ and $(i+1)\bmod
n=i\bmod n+1$.
\end{theorem}

\begin{theorem}[Quotient and remainder away from a boundary]
\lean{Nat.succ_div_mod_of_ne}
\label{Nat.succ_div_mod_of_ne}
\leanok
\uses{Nat.succ_div_mod_of_lt}
\difficulty{1}
For natural numbers $n$ and $i$, suppose $n>0$ and $i\bmod n+1\ne n$. Then $(i+1)/n=i/n$
and $(i+1)\bmod n=i\bmod n+1$.
\end{theorem}

\begin{theorem}[Quotient and remainder at a boundary]
\lean{Nat.succ_div_mod_of_eq}
\label{Nat.succ_div_mod_of_eq}
\leanok
\uses{Nat.mul_add_div_of_lt}
\difficulty{4}
For natural numbers $n$ and $i$, if $i$ divided by $n$ has remainder satisfying $i\bmod
n+1=n$, then $(i+1)/n=i/n+1$ and $(i+1)\bmod n=0$.
\end{theorem}

\begin{theorem}[Bound on the natural remainder]
\lean{Int.toNat_emod_lt}
\label{Int.toNat_emod_lt}
\leanok
\difficulty{3}
Let $M$ be a positive natural number and let $x$ be an integer. The natural-number
remainder of $x$ modulo $M$ is strictly less than $M$.
\end{theorem}

\begin{theorem}[Integer remainder is a natural number]
\lean{Int.natCast_toNat_emod}
\label{Int.natCast_toNat_emod}
\leanok
\difficulty{1}
Let $M$ be a positive natural number and let $x$ be an integer. The natural-number cast
of the remainder of $x$ modulo $M$, viewed again as an integer, equals that remainder.
\end{theorem}
